\subsection{Evaluation}

We evaluate soundtrack generation mainly on audio quality, audio-video alignment, and audio-text alignment.
We prioritize alignment to video over alignment to text, because text is used as a supplement and may not capture all the details in the video.
Moreover, text input is not presented to the viewers in the final output.
\cref{tab:audio_metric} summarizes the metrics. Correlations between subjective and objective metrics are studied in~\Cref{sec:app_audio_metric_corr}.


\begin{table}[h]
    \centering
    \adjustbox{max width=\textwidth}{%
    \begin{tabular}{l ll}
        \toprule
        &  \textbf{Subjective} & \textbf{Objective}\\
        \midrule
        \multirow{3}{*}{Audio quality}
        & Overall quality (Ovr.) & \multirow{3}{*}{Audio quality score (AQual)} \\
        & Naturalness (Nat.) \\
        & Professionalness (Pro.) \\
        \midrule
        \multirow{4}{*}{Video alignment}
        & Diegetic sound correctness (Corr.)& \multirow{4}{*}{\IB score (\IBShort)} \\
        & Diegetic sound synchronization (Sync.)\\
        & Non-diegetic music mood alignment (Mood) \\
        & Non-diegetic music motion/scene alignment (Action) \\
        \midrule
        \multirow{2}{*}{Text alignment}
        & Precision & \multirow{2}{*}{CLAP score} \\
        & Recall \\
        \bottomrule
    \end{tabular}
    }
    \caption{\textbf{\OursAudio evaluation metrics}. Audio quality measures the standalone quality of the generated audio, while video and text alignment measure how well the generated audio aligns with the input video and text respectively.}
    \label{tab:audio_metric}
\end{table}

\subsubsection{Metrics\label{sec:audiobox_metrics}}
\noindent\textbf{Audio quality.}
We aim to evaluate how \textit{natural} (free of artifacts) and \textit{professional} (\eg, volume balance, crispness) an audio sample sounds.
For subjective tests, a pairwise protocol is adopted which asks raters to choose which audio has better overall quality and on those two axes,
where the pair of audio samples are generated conditioned on the same video and text prompts when applicable.
We report the Net Win Rate (NWT), defined as ``win\% - lose\%'' for pairwise comparisons.
NWT ranges from $-100\%$  to $100\%$. Details are described at the end of the section.

For objective metric, the audio quality score (AQual) predicted by the model described in~\cref{sec:audio_cap} is used as an automatic metric.
We note that the model tends to assign higher scores to samples with music;
hence the metric should not be used when comparing samples with music and those without music.

Note that we do not adopt Frech\'et audio distance (FAD)~\citep{kilgour2018fr} or KL-divergence (KLD) metrics that are often reported in text-to-audio generation~\citep{liu2023audioldm, vyas2023audiobox}
because these metrics are not applicable to generated videos that do not have corresponding audio, which we mainly evaluate on in this paper.


\noindent\textbf{Video alignment.}
This measures how well the audio is aligned with the video. %
For diegetic sound, we measure \textit{correctness} and \textit{synchronization}.
Correctness reflects whether the \textit{right type} of audio is generated with respect to the scene and the objects in the video (\eg, dog barking versus cat meowing).
Synchronicity on the other hand focuses on whether the audio is generated at the \textit{right time} matching the motions in the video.
For non-diegetic background music, we measure how well it supports the \textit{mood} of the scene and how well the \textit{score synchronizes} with the on-screen actions and scene changes (\ie, action scoring).
Similarly, a pairwise protocol is adopted for subjective tests and net win rate is reported.

For automatic metrics, the \IB (\IBShort) score is used for measuring the alignment between video and diegetic sounds, as used in \citet{mei2023foleygen}.
As mentioned earlier, when non-diegetic music is present in the video,
the score usually decreases regardless whether the mood and the score matches because the model is trained on mostly diegetic data without non-diegetic music.


\noindent\textbf{Text alignment.}
Finally, we measure \textit{precision} (percentage of generated audio events that are in the text caption) and \textit{recall} (percentage of audio events from the caption that are generated in the audio).
We note that we focus more on recall than on precision, as the caption might not include all the acoustic events that are supposed to be heard in a video.
In terms of the subjective tests, raters are asked to rate on a scale from 1 to 5 for precision and recall.
We adopt the standalone protocol instead of a pairwise one because text alignment is more objective and is easier to rate in absolute scale.
For objective test, we measure this with CLAP score~\citep{wu2023clap},
which is commonly used for text-to-audio generation and does not distinguish hallucination and missing errors.

\noindent\textbf{Computing net win rate}
The reported subjective preference metric is Net Win Rate.
For a given model pair, NWT is computed as follows: each item is evaluated by three raters.
For each item evaluated, we take the mean of the preference between model $A$ and model $B$ (+1 if the model $A$ is preferred, 0 if a tie and -1 if the model $B$ is preferred).
These are the \textit{consensus} scores for each item.
We then average these consensus scores across all items to obtain a net win rate of $A$ (this is the expected fraction of items where model $A$ is preferred minus the fraction of items where model $B$ is preferred).
To obtain 95\% confidence intervals around Net Win Rate, we bootstrap resample the item-level consensus scores 1,000 times,
compute Net Win Rate for each, and take difference between the 2.5\%-ile and 97.5\%-ile of the Net Win Rate as the 95\% confidence interval.
The net win rate of $A$ \vs $B$ ranges from -100\% to 100\%.

\subsubsection{Audio Generation Benchmarks}
\label{subsubsec:audio_gen_bench}
To thoroughly evaluate audio generation, we consider multiple existing video sources including both real and generated videos,
and propose to release a benchmark \OursAudioBench\footnote{\url{https://github.com/facebookresearch/MovieGenBench}} which contains high quality videos generated by \OursVideo that cover a wide spectrum of audio events, and human reviewed sound and music captions for those videos.
In order to enable fair comparison to \OursAudio by future work, we also release non cherry picked generated audio from \OursAudio on \OursAudioBench.

We group videos into two categories: \textit{single-shot} and \textit{multi-shot}.
Single-shot videos are available in larger quantities and cover a wider spectrum of sound effects,
which are suitable for testing robustness and generalization.
Multi-shot videos, extracted from short films, contain scene transitions and have stronger sentiment and deeper narratives than single-shot videos.
Hence, they are suitable for evaluating video-music alignment and sound design perspectives,
such as when music enters, how music evolves with the story and aligns with the cuts, and whether music and sound effects are mixed harmonically.
We describe the composition of single-shot and multi-shot benchmarks next and \cref{tab:audio_eval_set} provides a summary.

\begin{table}[h]
    \centering
    \adjustbox{max width=\textwidth}{%
    \begin{tabular}{cc cc cc}
    \toprule
    Benchmark & Video sources & Type & Shot & Num. of samples &  Max length \\
    \midrule
    \OursAudioBenchSReal  & VGGSound            & Real & Mostly single  & 51    & 10s \\
    \OursAudioBenchSGen   & \RunwayGen, \Sora   & Gen  & Single         & 151   & 10s \\
    \OursAudioBench       & \OursVideo          & Gen  & Single         & 527   & 10s \\
    \OursAudioBenchMGen   & \OursVideo, \Sora   & Gen  & Multi          & 107   & 15s \\
    \bottomrule
    \end{tabular}
    }
\caption{\textbf{\OursAudio evaluation datasets}. We use "SFX" or "SFX+music" suffix to indicate what version of text captions is used.}
\label{tab:audio_eval_set}
\end{table}

\noindent\textbf{Single-shot.}
This includes VGGSound~\citep{vggsound}, \Sora~\citep{sora}, \RunwayGen~\citep{gen3}, and our proposed \OursAudioBench.
\begin{itemize}
    \item VGGSound contains real videos and is widely used for training and evaluating video-to-audio generation models~\citep{mei2023foleygen, luo2024diff, xing2024seeing}.
    However, we discovered that there are many duplicates or near duplicates (\eg, with added static watermark or text) between the training and evaluation split,
    and some testing videos are static.
    We perform deduplication based on video embeddings, and manually review test sets to select 51 samples that are not static, do not contain diegetic speech, and mostly have motion synchronized sounds.
    \item \Sora has been used to demonstrate video-to-audio generation for generated videos~\citep{xing2024seeing,mei2023foleygen}, but the number of available videos is small and of limited domain.
    Hence, it is not suitable for being used alone as a benchmark.
    We review those used in text-to-video comparison (\Cref{appendix_prior_work_details}) and selected 43 samples for audio generation evaluation.
    \item \OursAudioBench is the new benchmark dataset we create using \OursVideo.
    It includes 527 videos and is designed to cover various ambient environments (\eg, indoor, urban, nature, transportation) and sound effects (\eg, human, animal, objects).
    This is the first large scale synthetic benchmark for evaluating video-to-audio generation.
    To create this benchmark, we first define an ontology with 36 audio categories and audio concepts for each category (\eg, expression $\rightarrow$ \{cry, laugh, yell\}), with a total of 434 concepts.
    Llama3 is next used to propose video prompts for each audio concept,
    and \OursVideo is used to generate videos given these prompts.
    We next review the generated videos to exclude those with artifacts that would severely impact the judgement of whether an audio fits the video, resulting in the final 527 videos.
    \item \RunwayGen contains 108 synthetic videos.
    It is created with a similar process as \OursAudioBench but with a subset of prompts.
    The goal is to include synthetic videos from different models that may contain different types of artifacts and artistic styles, in order to test the robustness of video-to-audio generation models.
\end{itemize}

We group them into three sets based on video properties:
``\OursAudioBenchSReal'' is the real single-shot videos which includes VGGSound;
``\OursAudioBenchSGen'' is the generated single-shot videos from prior video generation models which include \Sora and \RunwayGen;
``\OursAudioBench'' is the new generated single-shot benchmark we create, which we hope will facilitate future work for thorough text and video-to-audio generation benchmarking.

\noindent\textbf{Multi-shot.}
We source 26 short films generated by \Sora and by \OursVideo to create this set.
These videos are 30 second to 2 minute long, and are composed of multiple related shots.
To compare our model with baseline methods, many of which have length limitations and do not support audio extension, we chunk these videos into 15-second segments and discard last segments shorter than 10 seconds, which results in 107 segments in total.
The combined set is referred to as \OursAudioBenchMGen.

\noindent\textbf{Text prompt creation for SFX and SFX+music generation.}
For all the video samples, we use Llama3~\citep{llama3} to propose 5 sound and music captions for each video given its video caption,
and manually select the best sound and music caption for each video.
To generate non-diegetic music and sound effects jointly,
we set the prompt to ``This audio contains music with a 0.90 likelihood.'' for the music presence part of the text caption (see \cref{tab:audio_caption}).
In contrast, the likelihood is set to 0.01 if music is undesired.
For baseline models that take text prompt as input, we use sound caption as input when generating sound effects only,
and concatenate sound and music caption when generating sound effect and music jointly.
